% GENERATED FILE. Edit canonical lessons, then run npm run generate:books.
% canonical-source-hash: fnv1a64:e2bb18c3e9f0f4ab
% canonical-lessons: ML-C363-jvala, ML-C363-nasanastam, ML-C363-alikkattuka, ML-C363-pottitterikkuka, ML-C363-takaruka

\chapter{Flames and Damage}
\label{ch:ml-flames-and-damage}
\hlchaptermodality{363}

\begin{chapteropening}
\textbf{By the end of this chapter:} \emph{I can say a flame, damage, to blaze up, to explode and to collapse.}

Complete the last lesson of chapter 363: I can say a flame, damage, to blaze up, to explode and to collapse.
\end{chapteropening}

\section[\texorpdfstring{\ml{ജ്വാല} (jvāla)}{ജ്വാല (jvāla)}]{\ml{ജ്വാല} (jvāla) — a flame}

\label{lesson:ML-C363-jvala}



\begin{quote}
Before the new one: say the Malayalam for to catch fire, then the Malayalam for to burn down.
\end{quote}

\subsection*{You'll want to know: \ml{ജ്വാല}}
\textbf{\ml{ജ്വാല}} — \emph{jvāla} — \textquotedblleft{}a flame\textquotedblright{}.

\begin{quote}
\ml{മെഴുകുതിരിയുടെ} \ml{ജ്വാല} \ml{കാറ്റിൽ} \ml{ഇളകി}.
\end{quote}

\begin{quote}
\emph{meḻukutiriyuṭe jvāla kāṟṟil iḷaki.}
\end{quote}

\begin{quote}
The candle's flame shook in the wind.
\end{quote}

A literary and news word; plural \ml{ജ്വാലകൾ}.

\begin{cousinweb}
\textbf{\ml{ജ്വാല}} (\emph{jvāla}) is borrowed from Sanskrit \emph{jvālā}, \textquotedblleft{}a flame, a blaze\textquotedblright{} (Monier-Williams), from the verb \emph{jval}, \textquotedblleft{}to burn brightly\textquotedblright{}.
\end{cousinweb}

\begin{grammarlens}[title={What you've built}]
One more word for telling what happened.
\end{grammarlens}

\subsection*{Your turn}
\begin{itemize}
\raggedright
  \item \emph{Say it:} \emph{jvāla}
  \item \emph{Say it:} \emph{jvāla}, once more
  \item \emph{Say it:} \emph{meḻukutiriyuṭe jvāla kāṟṟil iḷaki.} — The candle's flame shook in the wind.
  \item \emph{Recall:} say the Malayalam for to catch fire, then the Malayalam for to burn down, then say \emph{jvāla} again
\end{itemize}

\begin{tcolorbox}[breakable,colback=teal!4,colframe=teal!35!black,title={Before you move on}]
What does \ml{ജ്വാല} mean? (\textquotedblleft{}A flame\textquotedblright{}.) Say it once more.
\end{tcolorbox}
\section[\texorpdfstring{\ml{നാശനഷ്ടം} (nāśanaṣṭaṁ)}{നാശനഷ്ടം (nāśanaṣṭaṁ)}]{\ml{നാശനഷ്ടം} (nāśanaṣṭaṁ) — damage; losses}

\label{lesson:ML-C363-nasanastam}



\begin{quote}
Before the new one: say the Malayalam for to burn down, then the Malayalam for a flame.
\end{quote}

\subsection*{You'll want to know: \ml{നാശനഷ്ടം}}
\textbf{\ml{നാശനഷ്ടം}} — \emph{nāśanaṣṭaṁ} — \textquotedblleft{}damage; losses\textquotedblright{}.

\begin{quote}
\ml{വെള്ളപ്പൊക്കത്തിൽ} \ml{ഗ്രാമത്തിന്} \ml{വലിയ} \ml{നാശനഷ്ടം} \ml{ഉണ്ടായി}.
\end{quote}

\begin{quote}
\emph{veḷḷappokkattil grāmattinŭ valiya nāśanaṣṭaṁ uṇṭāyi.}
\end{quote}

\begin{quote}
The village suffered great damage in the flood.
\end{quote}

Two words that say almost the same, destruction and loss: the whole cost of a disaster.

\begin{cousinweb}
\textbf{\ml{നാശനഷ്ടം}} (\emph{nāśanaṣṭaṁ}) is built from two Sanskrit words: \emph{nāśa}, \textquotedblleft{}destruction, ruin\textquotedblright{} (CDIAL 7084), and \emph{naṣṭa}, \textquotedblleft{}lost, destroyed\textquotedblright{} (Monier-Williams), the source of \ml{നഷ്ടം} (\emph{naṣṭaṁ}), \textquotedblleft{}a loss\textquotedblright{}.
\end{cousinweb}

\begin{grammarlens}[title={What you've built}]
One more word for telling what happened.
\end{grammarlens}

\subsection*{Your turn}
\begin{itemize}
\raggedright
  \item \emph{Say it:} \emph{nāśanaṣṭaṁ}
  \item \emph{Say it:} \emph{nāśanaṣṭaṁ}, once more
  \item \emph{Say it:} \emph{veḷḷappokkattil grāmattinŭ valiya nāśanaṣṭaṁ uṇṭāyi.} — The village suffered great damage in the flood.
  \item \emph{Recall:} say the Malayalam for to burn down, then the Malayalam for a flame, then say \emph{nāśanaṣṭaṁ} again
\end{itemize}

\begin{tcolorbox}[breakable,colback=teal!4,colframe=teal!35!black,title={Before you move on}]
What does \ml{നാശനഷ്ടം} mean? (\textquotedblleft{}Damage; losses\textquotedblright{}.) Say it once more.
\end{tcolorbox}
\section[\texorpdfstring{\ml{ആളിക്കത്തുക} (āḷikkattuka)}{ആളിക്കത്തുക (āḷikkattuka)}]{\ml{ആളിക്കത്തുക} (āḷikkattuka) — to blaze up}

\label{lesson:ML-C363-alikkattuka}



\begin{quote}
Before the new one: say the Malayalam for a flame, then the Malayalam for damage.
\end{quote}

\subsection*{You'll want to know: \ml{ആളിക്കത്തുക}}
\textbf{\ml{ആളിക്കത്തുക}} — \emph{āḷikkattuka} — \textquotedblleft{}to blaze up\textquotedblright{}.

\begin{quote}
\ml{കാറ്റ്} \ml{വീശിയപ്പോൾ} \ml{തീ} \ml{ആളിക്കത്തി}.
\end{quote}

\begin{quote}
\emph{kāṟṟŭ vīśiyappōḷ tī āḷikkatti.}
\end{quote}

\begin{quote}
When the wind blew, the fire blazed up.
\end{quote}

Past: \ml{ആളിക്കത്തി}, \emph{āḷikkatti}.

Said of a fire that suddenly grows tall and fierce.

\begin{cousinweb}
\textbf{\ml{ആളിക്കത്തുക}} (\emph{āḷikkattuka}) joins \ml{ആളി}, \textquotedblleft{}flaring up\textquotedblright{}, and the native verb \ml{കത്തുക} (\emph{kattuka}), \textquotedblleft{}to burn\textquotedblright{} (DEDR 1207). The origin of the first part is not settled.
\end{cousinweb}

\begin{grammarlens}[title={What you've built}]
One more word for telling what happened.
\end{grammarlens}

\subsection*{Your turn}
\begin{itemize}
\raggedright
  \item \emph{Say it:} \emph{āḷikkattuka}
  \item \emph{Say it:} \emph{āḷikkattuka}, once more
  \item \emph{Say it:} \emph{kāṟṟŭ vīśiyappōḷ tī āḷikkatti.} — When the wind blew, the fire blazed up.
  \item \emph{Recall:} say the Malayalam for a flame, then the Malayalam for damage, then say \emph{āḷikkattuka} again
\end{itemize}

\begin{tcolorbox}[breakable,colback=teal!4,colframe=teal!35!black,title={Before you move on}]
What does \ml{ആളിക്കത്തുക} mean? (\textquotedblleft{}To blaze up\textquotedblright{}.) Say it once more.
\end{tcolorbox}
\section[\texorpdfstring{\ml{പൊട്ടിത്തെറിക്കുക} (poṭṭitteṟikkuka)}{പൊട്ടിത്തെറിക്കുക (poṭṭitteṟikkuka)}]{\ml{പൊട്ടിത്തെറിക്കുക} (poṭṭitteṟikkuka) — to explode}

\label{lesson:ML-C363-pottitterikkuka}



\begin{quote}
Before the new one: say the Malayalam for damage, then the Malayalam for to blaze up.
\end{quote}

\subsection*{You'll want to know: \ml{പൊട്ടിത്തെറിക്കുക}}
\textbf{\ml{പൊട്ടിത്തെറിക്കുക}} — \emph{poṭṭitteṟikkuka} — \textquotedblleft{}to explode\textquotedblright{}.

\begin{quote}
\ml{ഫാക്ടറിയിൽ} \ml{ഒരു} \ml{യന്ത്രം} \ml{പൊട്ടിത്തെറിച്ചു}.
\end{quote}

\begin{quote}
\emph{phākṭaṟiyil oru yantraṁ poṭṭitteṟiccu.}
\end{quote}

\begin{quote}
A machine exploded in the factory.
\end{quote}

Past: \ml{പൊട്ടിത്തെറിച്ചു}, \emph{poṭṭitteṟiccu}.

Also of a person: to explode with anger.

\begin{cousinweb}
\textbf{\ml{പൊട്ടിത്തെറിക്കുക}} (\emph{poṭṭitteṟikkuka}) is a native Dravidian compound of \ml{പൊട്ടുക} (\emph{poṭṭuka}), \textquotedblleft{}to burst, to crack\textquotedblright{} (DEDR 4490), and \ml{തെറിക്കുക} (\emph{teṟikkuka}), \textquotedblleft{}to fly off, to fly in pieces\textquotedblright{} (DEDR 3438): to burst so that pieces fly.
\end{cousinweb}

\begin{grammarlens}[title={What you've built}]
One more word for telling what happened.
\end{grammarlens}

\subsection*{Your turn}
\begin{itemize}
\raggedright
  \item \emph{Say it:} \emph{poṭṭitteṟikkuka}
  \item \emph{Say it:} \emph{poṭṭitteṟikkuka}, once more
  \item \emph{Say it:} \emph{phākṭaṟiyil oru yantraṁ poṭṭitteṟiccu.} — A machine exploded in the factory.
  \item \emph{Recall:} say the Malayalam for damage, then the Malayalam for to blaze up, then say \emph{poṭṭitteṟikkuka} again
\end{itemize}

\begin{tcolorbox}[breakable,colback=teal!4,colframe=teal!35!black,title={Before you move on}]
What does \ml{പൊട്ടിത്തെറിക്കുക} mean? (\textquotedblleft{}To explode\textquotedblright{}.) Say it once more.
\end{tcolorbox}
\section[\texorpdfstring{\ml{തകരുക} (takaruka)}{തകരുക (takaruka)}]{\ml{തകരുക} (takaruka) — to collapse; to be smashed}

\label{lesson:ML-C363-takaruka}



\begin{quote}
Before the new one: say the Malayalam for to blaze up, then the Malayalam for to explode.
\end{quote}

\subsection*{You'll want to know: \ml{തകരുക}}
\textbf{\ml{തകരുക}} — \emph{takaruka} — \textquotedblleft{}to collapse; to be smashed\textquotedblright{}.

\begin{quote}
\ml{ശക്തമായ} \ml{കാറ്റിൽ} \ml{പഴയ} \ml{പാലം} \ml{തകർന്നു}.
\end{quote}

\begin{quote}
\emph{śaktamāya kāṟṟil paḻaya pālaṁ takarnnu.}
\end{quote}

\begin{quote}
The old bridge collapsed in the strong wind.
\end{quote}

Past: \ml{തകർന്നു}, \emph{takarnnu}.

Also of hopes and plans: \ml{എന്റെ} \ml{സ്വപ്നം} \ml{തകർന്നു}, my dream was shattered.

\begin{cousinweb}
\textbf{\ml{തകരുക}} (\emph{takaruka}) is native Dravidian: DEDR 2999 lists Malayalam \emph{takaruka}, \textquotedblleft{}to be crushed, smashed\textquotedblright{}, beside Tamil \emph{takar}, \textquotedblleft{}to be broken to pieces\textquotedblright{}.
\end{cousinweb}

\begin{grammarlens}[title={What you've built}]
One more word for telling what happened.
\end{grammarlens}

\subsection*{Your turn}
\begin{itemize}
\raggedright
  \item \emph{Say it:} \emph{takaruka}
  \item \emph{Say it:} \emph{takaruka}, once more
  \item \emph{Say it:} \emph{śaktamāya kāṟṟil paḻaya pālaṁ takarnnu.} — The old bridge collapsed in the strong wind.
  \item \emph{Recall:} say the Malayalam for to blaze up, then the Malayalam for to explode, then say \emph{takaruka} again
\end{itemize}

\begin{tcolorbox}[breakable,colback=teal!4,colframe=teal!35!black,title={Before you move on}]
What does \ml{തകരുക} mean? (\textquotedblleft{}To collapse; to be smashed\textquotedblright{}.) Say it once more.
\end{tcolorbox}
